\uuid{vDlr}
\titre{Variables continues}
\auteur{worms}
\organisation{exo7}
\licence{CC BY-NC-SA 4.0}
\datecreate{2017-08-28}
\niveau{L2}
\module{Probabilité et statistique}
\chapitre{Probabilité continue}
\sousChapitre{Densité de probabilité}
\difficulte{}
\theme{Variables continues}
\exo7qcmid{worms.worms.042}

\contenu{
\begin{question}
\qclassification{261.01}
On considère que la quantité de pain (en centaines de kg) qu'une boulangerie vend en une journée est une variable aléatoire de loi de densité définie par
\[
 f_X(x) = \left\{  \begin{array}[c]{cl} x & \mbox{ si $0\leq x\leq 1$} \\ 2-x & \mbox{ si $1\leq x\leq 2$} \\ 0 & \mbox{ sinon}    \end{array} \right.
\]
On introduit les événements $A$="la boulangerie vendra demain au moins 100kg de pain" et $B$="la boulangerie vendra demain entre 50 et 150kg de pain". Les événements $A$ et $B$ sont ils indépendants ?
\begin{answers}
\good{Oui.}
\bad{Non.}
\end{answers}
\begin{explanations}
On a $A = \{ X \geq 1 \}$, $B = \{ 0.5 \leq X \leq 1.5 \}$, $A \cap B = \{ 1 \leq X \leq 1.5 \}$ ; il suffit alors de déterminer la fonction de répartition $F$ de $X$, et de vérifier que $\mathbb{P}( A \cap B ) = \mathbb{P}(A)  \mathbb{P}(B)$, c'est-à-dire que $F(1.5) - F(1) = (1-F(1)) (F(1.5) - F(0.5))$. On trouve alors qu'il y a bien égalité.
\end{explanations}
\end{question}
}
